%% hopfion-framework/gravity/cosmology.tex
%% Sources: main_paper7.tex
%% Inflation, n_s, r, P_s, N_e, reheating from condensate
%% Auto-generated from project files. Edit source papers to update.


%════════════════════════════════════════════════════════════════════
% Inflation and CMB Observables — Paper VII
%════════════════════════════════════════════════════════════════════
The Hopf-spoke formula of Paper~IV, via the Verlinde cancellation,
establishes that the ratio of the electron mass to the CMB
temperature is a dimensionless number determined entirely by the
topology and WZW algebra of the $Q=2$ icosahedral condensate:
\begin{equation}
  \frac{m_e}{\TCMB}
  = \left(\frac{\pi^2}{150}\right)^{\!\!1/4}
    \cdot\vp^{2Q}\cdot
    \exp\!\left(\frac{1600\pi - 3}{400}\right)
  \approx 2.18\times10^9
  \quad(k_B=1),
  \label{P7:eq:me_TCMB_ratio}
\end{equation}
with the right-hand side involving only $\pi$, $\vp$, $Q=10$, and the
WZW integers $k=3$, $c=9/5$; no Standard Model parameter enters. Paper
IV's Verlinde-cancellation theorem sharpens this to
\begin{equation}
  \frac{m_e}{\TCMB}
  \;=\; \left(\frac{\pi^2}{150}\right)^{\!1/4}
        \cdot\vp^{20}\cdot
        \exp\!\left(\frac{1600\pi-3}{400}\right)
        \cdot\bigl(1 + O(R_0^{-2})\bigr),
  \label{P7:eq:me_TCMB_ratio_corrected}
\end{equation}
a dimensionless geometric invariant of the $Q=2$ icosahedral
condensate accurate to $0.22\%$ ($O(R_0^{-2})$ thick-torus
correction). The framework previously required two experimental
inputs, $\TCMB$ and $m_e$; closing the $O(R_0^{-2})$ residual would
reduce this to one, and is closed in Paper~XII.

\begin{proposition}[Scale modulus: $\TCMB$ is not determined by the condensate dynamics]
\label{P7:prop:scale_modulus}
The density-feedback Faddeev--Niemi condensate does not fix the absolute
energy scale $\Lcond$.  $\TCMB$ is therefore the genuine single
experimental input of the framework not derivable from the condensate
dynamics alone.
\end{proposition}
\status{published}
\source{P7:prop:scale_modulus}
\begin{proof}
The density-feedback functional $\Efb[f;\beta]=\Kfb(\beta)\cdot\Jfour[f]$
is invariant under the joint rescaling
\begin{equation}
  f(\mathbf{x}) \;\to\; f(s\mathbf{x}),
  \qquad
  \beta \;\to\; \beta/s^2,
  \qquad
  \Lcond \;\to\; \Lcond/s,
  \label{P7:eq:scale_rescaling}
\end{equation}
for any $s>0$, since $\Jfour\sim\mathrm{length}^{-1}$,
$\Kfb\sim\mathrm{length}^{+1}$, so $\Efb=\Kfb\cdot\Jfour$ is
scale-invariant under $\mathbf{x}\to s\mathbf{x}$: no dynamical
equation fixes the absolute scale $\Lcond$, hence $\TCMB$ must be
supplied experimentally.
\end{proof}
\status{published}
\source{P7:prop:scale_modulus}

\begin{remark}[Invariance of $m_e/\TCMB$]
\label{P7:rem:me_tcmb_invariance}
Proposition~\ref{P7:prop:scale_modulus} explains why the ratio
$m_e/\TCMB$ (equation~\eqref{P7:eq:me_TCMB_ratio}) is a
\emph{dimensionless geometric invariant}: it is exactly the
combination $m_e/\Lcond \times \Lcond/\TCMB$ that is preserved by
the scale symmetry~\eqref{P7:eq:scale_rescaling}.
The ratio encodes pure condensate geometry; the absolute scale
$\TCMB$ cancels.
\end{remark}
\status{published}
\source{P7:rem:me_tcmb_invariance}

\begin{theorem}[CMB power spectrum from the see-saw scale]
\label{P7:thm:Ps_seesaw}
If the inflationary potential energy scale is $V_0 = M_{\mathrm{inf}}^4$,
the slow-roll CMB power spectrum normalisation for $N_e$ e-folds is
\begin{equation}
  \mathcal{P}_s
  = \frac{N_e^2}{12\pi^2}\left(\frac{M_{\mathrm{inf}}}{\MPl}\right)^4
  = \frac{N_e^2}{12\pi^2}\left(\frac{\LUV}{\MPl}\right)^8.
  \label{P7:eq:Ps_seesaw}
\end{equation}
For $N_e=60$ (instantaneous reheating): $\mathcal{P}_s = 2.82\times10^{-9}$,
a factor of $1.34$ above the observed $2.10\times10^{-9}$.
\end{theorem}
\status{published}
\source{P7:thm:Ps_seesaw}
\begin{proof}
Standard slow-roll result $\mathcal{P}_s=H^2/(8\pi^2\epsilon\MPl^2)$
with $H^2\approx V_0/(3\MPl^2)$ and $\epsilon=2/N_e^2$ (Starobinsky
plateau) gives $\mathcal{P}_s=N_e^2V_0/(12\pi^2\MPl^4)$. Substituting
$V_0=M_{\mathrm{inf}}^4=(\LUV^2/\MPl)^4$ and $(\LUV/\MPl)^2=1/\vp^{12}$
(from~\eqref{P7:eq:MPl}):
\[
  \mathcal{P}_s
  = \frac{N_e^2}{12\pi^2}\,\frac{\LUV^8}{\MPl^8}
  \approx \frac{3600}{118.4}\times(0.0557)^8
  = 2.82\times10^{-9}.
\]
\end{proof}
\status{published}
\source{P7:thm:Ps_seesaw}

\begin{proposition}[Gravitational see-saw scale]
\label{P7:prop:seesaw}
The condensate UV scale $\LUV=\MPl/\vp^6$ and the Planck mass
determine a gravitational see-saw scale
\begin{equation}
  \boxed{M_{\mathrm{inf}} \;=\; \frac{\LUV^2}{\MPl}
  \;=\; \frac{\MPl}{\vp^{12}}
  \;\approx\; (0.0557)^2\,\MPl
  \;\approx\; 3.79\times10^{16}\,\mathrm{GeV}.}
  \label{P7:eq:Minf}
\end{equation}
This scale is structurally analogous to the neutrino see-saw
$m_\nu\sim v_{\mathrm{EW}}^2/M_{\mathrm{GUT}}$: two derived scales
combine to produce an intermediate scale requiring no new parameters.
\end{proposition}
\status{published}
\source{P7:prop:seesaw}
\begin{proof}
Substitute $\LUV=\MPl/\vp^6$ into $\LUV^2/\MPl$:
$M_{\mathrm{inf}}=\MPl/\vp^{12}=(0.0557)^2\,\MPl=3.10\times10^{-3}\,\MPl
=3.79\times10^{16}\,\mathrm{GeV}$ exactly.
\end{proof}
\status{published}
\source{P7:prop:seesaw}

\begin{proposition}[$O(R_0^{-2})$ correction to the see-saw scale]
\label{P7:prop:seesaw_correction}
The exact 3D virial (Lemma~3.2 of~\cite{Paper4}) implies that $\LUV$
receives a finite-torus correction
\begin{equation}
  \LUV(R_0) = \LUV(\infty)\cdot\sqrt{\frac{2R_0^2}{2R_0^2+1}},
  \label{P7:eq:LUV_correction}
\end{equation}
since $\LUV\sim 1/C^*$ and $C^{*2}$ shifts by the factor
$(2R_0^2+1)/(2R_0^2)$ from the winding virial correction.
At $R_0=3$ this gives $\LUV(3)=\LUV(\infty)\cdot\sqrt{18/19}$,
reducing $M_{\mathrm{inf}}^4$ by $(18/19)^4=0.8055$ and hence
\begin{equation}
  \mathcal{P}_s(R_0=3) = 2.82\times10^{-9}\times\frac{18^4}{19^4}
  = 2.27\times10^{-9}.
  \label{P7:eq:Ps_corrected}
\end{equation}
The residual factor after this correction is $2.27/2.10 = 1.08$.
\end{proposition}
\status{published}
\source{P7:prop:seesaw_correction}

\begin{proposition}[CMB normalisation from gravitational reheating]
\label{P7:prop:Treh}
The residual factor $1.08$ in $\mathcal{P}_s$ (Proposition~\ref{P7:prop:seesaw_correction})
is closed by gravitational reheating.
For a plateau inflaton of mass $m_\phi\sim M_{\mathrm{inf}}$ decaying through
Planck-suppressed dimension-5 operators, the reheating temperature satisfies
\begin{equation}
  T_{\mathrm{reh}} \;\sim\; \left(\frac{m_\phi^3}{\MPl^2}\right)^{1/2}
  \;\approx\; 10^7\text{--}10^9\,\mathrm{GeV},
  \label{P7:eq:Treh}
\end{equation}
with no free parameter (the coupling is $\sim1/\MPl$ by dimensional analysis).
Substituting into the e-fold formula~\eqref{P7:eq:Ne_window} with
$M_{\mathrm{inf}}=3.79\times10^{16}\,\mathrm{GeV}$:
\begin{equation}
  N_e \;=\; 57.6 - \underbrace{\ln\!\frac{M_{\mathrm{inf}}}{10^{16}\,\mathrm{GeV}}}_{1.33}
              - \underbrace{\frac{1}{3}\ln\!\frac{T_{\mathrm{reh}}}{10^{10}\,\mathrm{GeV}}}_{-2.4\text{ to }-3.4}
  \;\approx\; 57.4\text{--}58.1.
  \label{P7:eq:Ne_window}
\end{equation}
The exact normalisation $\mathcal{P}_s = 2.10\times10^{-9}$ is reached at
$N_e = 57.73$, which lies within~\eqref{P7:eq:Ne_window} for
$T_{\mathrm{reh}}\approx1.25\times10^{8}\,\mathrm{GeV}$,
a value consistent with gravitational reheating.
The Planck~2018 one-sigma band $\mathcal{P}_s\in[2.07,2.13]\times10^{-9}$
is satisfied for the broad window
\begin{equation}
  T_{\mathrm{reh}}\;\in\;[3.6\times10^{7},\,4.3\times10^{8}]\,\mathrm{GeV},
  \label{P7:eq:Treh_window}
\end{equation}
requiring no fine-tuning.
The corresponding spectral index and tensor-to-scalar ratio are
\begin{equation}
  \boxed{n_s = 1-\frac{2}{N_e} \approx 0.9654
  \quad(\text{within }0.12\sigma\text{ of Planck~2018}),
  \qquad
  r = \frac{12}{N_e^2} \approx 0.0036,}
  \label{P7:eq:inflation_final}
\end{equation}
with $\mathcal{P}_s = 2.10\times10^{-9}$ exact.
\end{proposition}
\status{published}
\source{P7:prop:Treh}

The kinetic term of the parent action is non-canonical. Defining the
canonical field
\begin{equation}
  \varphi_c(\rho) = \frac{2}{\vp^3\beta}\!\left[\sqrt{1+\beta\rho}-1\right],
  \label{P7:eq:phi_c}
\end{equation}
the Klein--Gordon equation on FLRW takes the standard form
$\ddot\varphi_c+3H\dot\varphi_c+V'(\varphi_c)=0$, where
\begin{equation}
  V(\varphi_c) = |\Lambda_0|\,e^{-\vp^{12}\varphi_c^2/4}
  \quad(\beta\rho\gg1,\;\text{early universe}),
  \label{P7:eq:Vinf}
\end{equation}
a Gaussian potential admitting a slow-roll plateau for $\varphi_c\ll2/\vp^6$.
\status{published}
\source{P7:sec:inflation}

\begin{remark}[CMB normalisation: robust spectral predictions]
\label{P7:rem:inflation_gap}
The tree-level inflationary potential is too steep for slow roll
($\epsilon_V\gg1$ at the Gaussian potential): the $\xNMC R$ branch
requires the non-minimal coupling
\begin{equation}
  \xNMC = \frac{\LUV^2\beta}{4(4\pi)^2},
  \label{P7:eq:xNMC}
\end{equation}
generated by the one-loop action.
The predictions~\eqref{P7:eq:inflation_predictions} for $n_s$ and $r$
are robust in the Starobinsky limit and independent of the overall
normalisation.
\end{remark}
\status{published}
\source{P7:rem:inflation_gap}

\begin{remark}[Reheating and the dark-matter initial conditions]
\label{P7:rem:reheating_IC}
Reheating after Higgs/Starobinsky inflation
(Proposition~\ref{P7:prop:Treh}, $T_{\rm reh}\approx1.25\times10^8\,\mathrm{GeV}$)
populates radiation and establishes the condensate --- the nematic
medium in which the director then orders. The dark-matter initial
conditions follow from the director sector. The cold, baryon-sourced
strain (Section~\ref{P7:sec:cold_director}) inherits the adiabatic
spectrum $\mathcal{P}_s=2.10\times10^{-9}$ (Proposition~\ref{P7:prop:Treh})
through the baryons it tracks, with no additional free parameter. The
splay self-organisation (Section~\ref{P7:sec:splay_instability}) adds a
\emph{cold} texture component, seeded at the orientational ordering
transition: an isocurvature contribution. Because the director is
ultralight ($m_\xi\sim H_0$), it is Hubble-frozen until $H\sim m_\xi$
and orders only at $z\sim0$; the texture therefore forms in the
present epoch, is absent at recombination, and carries negligible
CMB isocurvature, safely below the Planck bound. The same lateness
fixes what this texture component is: a late-time, galactic-scale
contribution, not a source of dark matter at recombination. Whether
the adiabatic baryon-sourced strain instead operates at recombination
depends on the epoch at which the nematic medium orders --- set by the
same condensate cutoff and not fixed here; in either case the strain
shifts no acoustic scale. The earlier proposal that the condensate is
itself the inflaton, giving $\varepsilon\approx\beta/\vp^6$, is
excluded by the Higgs/Starobinsky value $\varepsilon=2/N_e^2\approx0.001$.
\end{remark}
\status{published}
\source{P7:rem:reheating_IC}

\begin{corollary}[Tower level of the ultraviolet cutoff]
\label{P7:cor:luv_tower}
With $\LUV=\MPl/\vp^6$~\eqref{P7:eq:MPl}, the Planck mass and the
cutoff are locked to the tower spacing $\vp^6$:
\begin{equation}
  \LUV \;=\; \frac{\MPl}{\vp^6}\;\approx\;0.0557\,\MPl,
  \label{P7:eq:LUV_from_MPl}
\end{equation}
the gap $\vp^6=\lambda$ being the Bogomolny parameter that also fixes
$G_N$. Within the bosonic tower $M_n = \Lcond\,\vp^{2n}$ of
Paper~VI~\cite{Paper6}, the cutoff lies at the non-integer level
\begin{equation}
  \boxed{n_{\mathrm{UV}}
  \;=\; \frac{\ln(\LUV/\Lcond)}{2\ln\vp}
  \;\approx\; 73.6,}
  \label{P7:eq:nUV_formula}
\end{equation}
with $\Lcond = \TCMB(\pi^2/150)^{1/4}\approx1.19\times10^{-4}\,\mathrm{eV}$
the condensate ground state, while $\MPl$ lies three steps higher at
$n\approx76.6$. The $\vp^6$ gap between them is exact, but neither
absolute level is an integer: the cutoff is locked to the tower
spacing without coinciding with a discrete tower state.
\end{corollary}
\status{published}
\source{P7:cor:luv_tower}
\begin{proof}
From $\LUV=\MPl/\vp^6\approx0.0557\,\MPl\approx6.8\times10^{17}\,\mathrm{GeV}$:
\[
  n_{\mathrm{UV}}
  = \frac{\ln(\LUV/\Lcond)}{2\ln\vp}
  \approx \frac{\ln(6.8\times10^{17}\,\mathrm{GeV}\,/\,1.19\times10^{-13}\,\mathrm{GeV})}{2\times0.4812}
  \approx \frac{70.82}{0.9624}
  \approx 73.6,
\]
and $n(\MPl)=n_{\mathrm{UV}}+3\approx76.6$ since $\MPl=\vp^6\LUV$.
Since $73.6\notin\mathbb{Z}$, $\LUV$ does not coincide with any
discrete tower level, while the gap $n(\MPl)-n_{\mathrm{UV}}=3$
(i.e.\ $\vp^6$) is exact.
\end{proof}
\status{published}
\source{P7:cor:luv_tower}

\begin{remark}[Connection to Paper~VI]
\label{P7:rem:bogomolny_coincidence}
Corollary~\ref{P7:cor:luv_tower} is the result that Paper~VI~\cite{Paper6}
states as its final theorem (Theorem~7.1 therein). The $\vp^6$ gap
between $\MPl$ and $\LUV$~\eqref{P7:eq:MPl} is the Bogomolny parameter
$\lambda=\vp^6$ --- the same quantity that fixes $G_N$, the condensate
energy, and the $\vp^{-2}$ tower spacing --- so the cutoff sits three
tower steps below the Planck mass for a topological reason. (The
Seeley--DeWitt one-loop coefficient $4\pi\sqrt{2}\approx17.77$ happens
to lie within $1\%$ of $\vp^6\approx17.94$, but this loop--tower
proximity is a numerical coincidence and plays no role: the cutoff is
fixed by the tower relation $\LUV=\MPl/\vp^6$, and the one-loop term
only confirms that a gravitational sector is induced at the
condensate scale.) Together, Papers~VI and~VII close the
tower-to-gravity connection: Paper~VI establishes the $\vp$-tower from
first principles; Paper~VII locates the induced gravitational sector
within it.
\end{remark}
\status{published}
\source{P7:rem:bogomolny_coincidence}

\begin{equation}
  n_s = 1 - \frac{2}{N_e} = 0.9667,
  \qquad
  r = \frac{12}{N_e^2} = 0.0033.
  \label{P7:eq:inflation_predictions}
\end{equation}

